\section{Introduction}
This chapter presents the logical and structural design of GestureTalk. It describes how components interact, how data moves from the camera to speech, and how the main runtime objects are related. The design is shown through architecture, data-flow, journey, entity, use-case and sequence diagrams.

\section{Architecture Overview}
The system is a single Windows process. The main thread captures frames and classifies. A daemon thread owns the speech engine. Figure~\ref{fig:arch} shows the high-level path from the user hand to text and audio.

\begin{figure}[H]
\centering
\begin{adjustbox}{max width=\textwidth,center}
\begin{tikzpicture}[
  node distance=1.05cm and 0.45cm,
  box/.style={draw, thick, rounded corners, minimum width=2.15cm, minimum height=1.0cm, align=center, fill=white, font=\small},
  arr/.style={-Stealth, thick, black}
]
\node[box] (user) {User\\hand pose};
\node[box, right=of user] (cam) {Webcam\\OpenCV};
\node[box, right=of cam] (mp) {MediaPipe\\21 landmarks};
\node[box, right=of mp] (cls) {Rule\\classifier};
\node[box, below=of cls] (sm) {Smoothing\\+ confirm};
\node[box, left=of sm] (ui) {On-screen\\word};
\node[box, left=of ui] (tts) {SAPI\\speech};

\draw[arr] (user) -- (cam);
\draw[arr] (cam) -- (mp);
\draw[arr] (mp) -- (cls);
\draw[arr] (cls) -- (sm);
\draw[arr] (sm) -- (ui);
\draw[arr] (sm) -- (tts);
\draw[arr] (ui.west) -- ++(-0.2,0) |- (user.south);
\end{tikzpicture}

\end{adjustbox}
\caption{System architecture of GestureTalk}
\label{fig:arch}
\end{figure}

The capture layer uses OpenCV for the webcam, a horizontal flip and the display window \cite{opencv}. The perception layer uses MediaPipe Tasks HandLandmarker in VIDEO mode \cite{mediapipehl}. The classification layer applies geometry rules and named poses over a vocabulary of more than five hundred gestures. The temporal layer applies hysteresis and hold-to-confirm. The output layer draws the word on the frame and speaks it through interruptible Windows SAPI \cite{sapi}.

\section{DFD Level 0 (Context Diagram)}
\begin{figure}[H]
\centering
\begin{adjustbox}{max width=\textwidth,center}
\begin{tikzpicture}[
  node distance=1.45cm and 1.8cm,
  ext/.style={draw, thick, minimum width=2.4cm, minimum height=0.9cm, fill=white, align=center, font=\small},
  proc/.style={draw, thick, rounded corners, minimum width=4.0cm, minimum height=1.15cm, fill=white, align=center},
  arr/.style={-Stealth, thick, black}
]
\node[proc] (sys) {GestureTalk};
\node[ext, left=of sys] (user) {User};
\node[ext, above=of sys] (cam) {Webcam};
\node[ext, right=of sys] (spk) {Speakers};
\node[ext, below=of sys] (model) {hand\_landmarker.task};

\draw[arr] (user) -- node[above,font=\scriptsize]{pose} (sys);
\draw[arr] (sys) -- node[below,font=\scriptsize]{text} (user);
\draw[arr] (cam) -- node[right,font=\scriptsize]{frames} (sys);
\draw[arr] (sys) -- node[above,font=\scriptsize]{audio} (spk);
\draw[arr] (sys) -- (model);
\draw[arr] (model) -- (sys);
\end{tikzpicture}

\end{adjustbox}
\caption{DFD Level 0}
\label{fig:dfd0}
\end{figure}

At the context level, GestureTalk exchanges video with the webcam, landmark inference with the on-device model file, visual feedback with the user, and audio with the speakers. Internal processing is abstracted as a single system boundary.

\section{DFD Level 1}
\begin{figure}[H]
\centering
\begin{adjustbox}{max width=\textwidth,center}
\begin{tikzpicture}[
  node distance=0.75cm and 0.95cm,
  p/.style={draw, thick, rounded corners, minimum width=3.1cm, minimum height=0.8cm, align=center, fill=white, font=\small},
  d/.style={draw, thick, minimum width=2.5cm, minimum height=0.65cm, align=center, fill=white, font=\small},
  arr/.style={-Stealth, thick, black}
]
\node[p] (p1) {P1: Frame capture};
\node[p, below=of p1] (p2) {P2: Landmark extract};
\node[p, below=of p2] (p3) {P3: Classify pose};
\node[p, below=of p3] (p4) {P4: Confirm / speak};
\node[d, right=2.2cm of p1] (d1) {D1: Frame};
\node[d, below=of d1] (d2) {D2: Landmarks};
\node[d, below=of d2] (d3) {D3: Recent labels};
\node[d, below=of d3] (d4) {D4: Speech queue};

\draw[arr] (p1) -- (d1);
\draw[arr] (d1) -- (p2);
\draw[arr] (p2) -- (d2);
\draw[arr] (d2) -- (p3);
\draw[arr] (p3) -- (d3);
\draw[arr] (d3) -- (p4);
\draw[arr] (p4) -- (d4);
\end{tikzpicture}

\end{adjustbox}
\caption{DFD Level 1}
\label{fig:dfd1}
\end{figure}

Level~1 splits the system into capture, landmark extraction, classification, confirmation, display and speech. Each process reads a small set of stores: the current frame, the landmark buffer, recent labels, the stable word and the speech queue. Shared stores make flicker control and interruption explicit.

\section{DFD Level 2 (Classification and Confirmation)}
\begin{figure}[H]
\centering
\begin{adjustbox}{max width=\textwidth,center}
\begin{tikzpicture}[
  node distance=0.6cm,
  p/.style={draw, thick, rounded corners, minimum width=5.4cm, minimum height=0.68cm, align=center, fill=white, font=\small},
  d/.style={draw, thick, minimum width=5.4cm, minimum height=0.62cm, align=center, fill=white, font=\small},
  arr/.style={-Stealth, thick, black}
]
\node[p] (a) {P3.1 Compute finger ratios and orientation};
\node[p, below=of a] (b) {P3.2 Apply open/close hysteresis};
\node[p, below=of b] (c) {P3.3 Named rule or 500+ word vocabulary};
\node[p, below=of c] (d) {P3.4 Majority vote; else UNCERTAIN};
\node[p, below=of d] (e) {P3.5 Hold 0.1 s then speak once};
\node[d, below=of e] (s) {D3 Recent labels / D4 Speech queue};

\draw[arr] (a) -- (b);
\draw[arr] (b) -- (c);
\draw[arr] (c) -- (d);
\draw[arr] (d) -- (e);
\draw[arr] (e) -- (s);
\end{tikzpicture}

\end{adjustbox}
\caption{DFD Level 2}
\label{fig:dfd2}
\end{figure}

Level~2 expands classification into finger ratios, hysteresis, vocabulary lookup, majority vote and hold-to-confirm. Finger ratios are compared with separate open and close thresholds. Named rules such as thumbs-down NO and the OK pinch override the table. Disagreeing frames become uncertain and do not replace the last stable word.

\section{User Journey Diagram}
\begin{figure}[H]
\centering
\begin{adjustbox}{max width=\textwidth,center}
\begin{tikzpicture}[
  node distance=0.48cm,
  box/.style={draw, thick, rounded corners, minimum width=6.0cm, minimum height=0.65cm, align=center, fill=white, font=\small},
  arr/.style={-Stealth, thick, black}
]
\node[box] (a) {Start: python scripts/sign\_language\_app.py};
\node[box, below=of a] (b) {Webcam opens; mirrored view};
\node[box, below=of b] (c) {Show one hand; skeleton overlay};
\node[box, below=of c] (d) {Hold pose $\approx$ 0.1 s};
\node[box, below=of d] (e) {Word on screen + spoken once};
\node[box, below=of e] (f) {Change pose $\rightarrow$ interrupt and speak new word};
\node[box, below=of f] (g) {Remove hand $\geq$ 0.15 s $\rightarrow$ silence};
\node[box, below=of g] (h) {Press Q to quit};

\draw[arr] (a) -- (b);
\draw[arr] (b) -- (c);
\draw[arr] (c) -- (d);
\draw[arr] (d) -- (e);
\draw[arr] (e) -- (f);
\draw[arr] (f) -- (g);
\draw[arr] (g) -- (h);
\end{tikzpicture}

\end{adjustbox}
\caption{End-to-end user journey from application start to spoken word}
\label{fig:journey}
\end{figure}

The journey starts with launching \code{scripts/sign_language_app.py}. The webcam opens, the user shows one hand, holds a pose for about 0.1 seconds, sees and hears the word, may change pose, may remove the hand, and finally presses Q to quit.

\section{Entity Relationship Diagram}
There is no MySQL schema. Figure~\ref{fig:er} models runtime objects and their relationships.

\begin{figure}[H]
\centering
\begin{adjustbox}{max width=\textwidth,center}
\begin{tikzpicture}[
  node distance=1.2cm and 1.4cm,
  ent/.style={draw, thick, minimum width=3.0cm, minimum height=1.1cm, align=center, fill=white, font=\small},
  arr/.style={-Stealth, thick, black}
]
\node[ent] (h) {HandResult\\landmarks, side, score};
\node[ent, right=of h] (c) {GestureClassifier\\finger state, window};
\node[ent, below=of h] (g) {GestureState\\candidate, stable};
\node[ent, below=of c] (s) {SpeechController\\queue, generation};

\draw[arr] (h) -- node[above,font=\scriptsize]{feeds} (c);
\draw[arr] (c) -- node[right,font=\scriptsize]{detected word} (g);
\draw[arr] (g) -- node[above,font=\scriptsize]{say / stop} (s);
\end{tikzpicture}

\end{adjustbox}
\caption{Runtime entity relationship diagram}
\label{fig:er}
\end{figure}

One hand result exists per detected hand, with a maximum of one hand in the current demo. The classifier owns finger hysteresis. Gesture state owns confirmation. The speech controller owns the queue and generation counter. These objects live in memory only.

\section{Use Case Diagram}
\begin{figure}[H]
\centering
\begin{adjustbox}{max width=\textwidth,center}
\begin{tikzpicture}[
  node distance=0.7cm and 1.2cm,
  actor/.style={draw, thick, minimum width=2.4cm, minimum height=0.7cm, fill=white, align=center, font=\small},
  use/.style={draw, thick, ellipse, minimum width=4.2cm, minimum height=0.8cm, fill=white, align=center, font=\small},
  arr/.style={-Stealth, thick, black}
]
\node[actor] (u) {Presenter};
\node[actor, below=2.2cm of u] (d) {Developer};

\node[use, right=2.2cm of u] (u1) {Show pose / hear word};
\node[use, below=of u1] (u2) {Change pose / interrupt};
\node[use, below=of u2] (u3) {Remove hand / silence};
\node[use, below=of u3] (u4) {Toggle Left/Right (H)};
\node[use, below=of u4] (u5) {Run synthetic tests};

\draw[arr] (u) -- (u1);
\draw[arr] (u) -- (u2);
\draw[arr] (u) -- (u3);
\draw[arr] (u) -- (u4);
\draw[arr] (d) -- (u5);
\end{tikzpicture}

\end{adjustbox}
\caption{Use case diagram}
\label{fig:usecase}
\end{figure}

The presenter shows a pose, changes pose, removes the hand and may toggle left and right correction. The developer runs synthetic tests.

\section{Sequence Diagram (Stable Gesture to Speech)}
\begin{figure}[H]
\centering
\begin{adjustbox}{max width=\textwidth,center}
\begin{tikzpicture}[
  font=\footnotesize,
  actor/.style={draw, thick, minimum width=1.85cm, minimum height=0.5cm, align=center, fill=white},
  arr/.style={-Stealth, thick, black}
]
\node[actor] (w) {Webcam};
\node[actor, right=0.85cm of w] (t) {Tracker};
\node[actor, right=0.85cm of t] (c) {Classifier};
\node[actor, right=0.85cm of c] (g) {GestureState};
\node[actor, right=0.85cm of g] (s) {Speech};

\foreach \n in {w,t,c,g,s}
  \draw[dashed] (\n.south) -- ++(0,-3.7);

\draw[arr] ([yshift=-0.65cm]w.south) -- node[above,font=\scriptsize]{frame} ([yshift=-0.65cm]t.south);
\draw[arr] ([yshift=-1.35cm]t.south) -- node[above,font=\scriptsize]{21 pts} ([yshift=-1.35cm]c.south);
\draw[arr] ([yshift=-2.05cm]c.south) -- node[above,font=\scriptsize]{word} ([yshift=-2.05cm]g.south);
\draw[arr] ([yshift=-2.75cm]g.south) -- node[above,font=\scriptsize]{say()} ([yshift=-2.75cm]s.south);
\draw[arr] ([yshift=-3.4cm]s.south) -- node[below,font=\scriptsize]{audio} ++(-7.4,0);
\end{tikzpicture}

\end{adjustbox}
\caption{Sequence diagram for a confirmed gesture}
\label{fig:seq}
\end{figure}

The webcam loop sends a frame to the tracker, receives twenty-one points, asks the classifier for a word, waits for GestureState to confirm, and then asks the speech thread to speak. Audio returns to the user without blocking the next camera frame.

\section{Module Design}
The webcam loop, confirmation window, overlay and key handling live in \code{scripts/sign_language_app.py}. MediaPipe wrapping and skeleton drawing live in \code{src/signlang/hand_tracker.py}. Vocabulary, geometry and hysteresis live in \code{src/signlang/gestures.py}. Interruptible SAPI lives in \code{src/signlang/speech.py}. Synthetic landmark tests live in \code{tests/test_gestures.py}. The MediaPipe model file is stored at \code{models/hand_landmarker.task}.

\begin{table}[H]
\centering
\small
\caption{Source modules and responsibilities}
\label{tab:modules}
\begin{tabularx}{\textwidth}{@{}>{\raggedright\arraybackslash}p{3.6cm} L@{}}
\toprule
\textbf{Module} & \textbf{Responsibility} \\
\midrule
sign\_language\_app.py & Frame loop, confirm window, keys Q and H, overlay text \\
hand\_tracker.py & VIDEO-mode landmark detection and drawing \\
gestures.py & 500+ gesture vocabulary, named poses, handedness correction \\
speech.py & Background thread, generation counter, SAPI purge \\
test\_gestures.py & Synthetic coverage and noise protocol \\
check\_env.py & Python, packages, model, webcam and TTS checks \\
\bottomrule
\end{tabularx}
\end{table}
